% Generated draft: requires mathematical and source review.
\documentclass[11pt]{article}
\usepackage[margin=25mm]{geometry}
\usepackage{fontspec}
\setmainfont{Latin Modern Roman}
\usepackage{amsmath,amssymb,mathtools,unicode-math}
\setmathfont{Latin Modern Math}
\usepackage{enumitem,needspace,xurl,hyperref}
\setlength{\parindent}{0pt}
\setlength{\parskip}{5pt}
\setlength{\emergencystretch}{3em}
\newcommand{\N}{\mathbb{N}}
\newcommand{\Z}{\mathbb{Z}}
\newcommand{\Q}{\mathbb{Q}}
\newcommand{\R}{\mathbb{R}}
\newcommand{\C}{\mathbb{C}}
\newcommand{\F}{\mathbb{F}}
\newcommand{\A}{\mathbb{A}}
\newcommand{\PP}{\mathbb{P}}
\newcommand{\E}{\mathbb{E}}
\newcommand{\eps}{\varepsilon}
\newcommand{\dd}{\,\mathrm{d}}
\begin{document}
\section*{UnsolvedMath ID: 30000990}
\textbf{Source problem code:} catalogue entry
\subsection{Definitions and mathematical statement}
Let \(K\) be a filtered simplicial complex, and let \(P\) be its persistence pairing as intended in the report. In the usual homological setting a pair \((\sigma,\tau)\) consists of a birth simplex \(\sigma\) of dimension \(q\) and a death simplex \(\tau\) of dimension \(q+1\). Form the partial matching
\[
V_P=\{(\alpha,\beta)\in P:\alpha\text{ is a codimension-one face of }\beta\}.
\]
The face condition has codimension one under the usual dimension convention for persistence pairs.

Conjecture 2 in the source asks: for every persistence pair \((\sigma,\tau)\in P\) with \(\sigma\) not a face of \(\tau\), is there exactly one gradient path from \(\tau\) to \(\sigma\) in \(V_P\)? A customary combinatorial meaning of such a path is an alternating sequence
\[
\tau>\alpha_0<\beta_0>\alpha_1<\beta_1>\cdots<\beta_{r-1}>\alpha_r=\sigma,
\]
where \(\alpha_0\) is a \(q\)-dimensional face of \(\tau\), each \((\alpha_i,\beta_i)\in V_P\), and \(\alpha_{i+1}\ne\alpha_i\) is another \(q\)-dimensional face of \(\beta_i\). The notation \(>\) denotes the indicated face incidence. Under this convention the assertion is that the set of these paths has cardinality exactly one, including existence, not just a statement about a signed or mod-two count.

The filtration finiteness, simplex ordering and tie-breaking rules, coefficient field, choice of persistence pairing, and precise gradient-path convention are not given. They must be recovered before asserting the conjecture for every possible filtration or reduction procedure. The source calls \(V_P\) a discrete vector field but does not reproduce any acyclicity result needed for its use as a discrete gradient field.
\subsection{Short English statement}
For a persistence pair whose simplices are not incident, is there exactly one alternating gradient path from the death simplex to the birth simplex using the incident persistence pairs? The source’s pairing and path conventions require confirmation.
\subsection{Sources}
\begin{itemize}
\item Local UnsolvedMath registry, permanent ID \(30000990\), source code catalogue entry, ``Gradient-Path Uniqueness in Persistence Pairings.'' The supplied assessment identifies this as Conjecture 2; no individual proposer is supplied.
\item \emph{Computational Algebraic Topology} (2009), Oberwolfach Reports, DOI 10.4171/owr/2008/29, pp. 1603--1654: \url{https://doi.org/10.4171/owr/2008/29}. This is the sole supplied external reference.
\item Supplied literature triage, checked 2026-08-21: status ``open.'' It says the report explicitly presents the single-gradient-path assertion as Conjecture 2 and reports no located authoritative proof or counterexample. This dated assessment derives from supplied OWR curation; time-sensitive status and cited claims require independent verification.
\end{itemize}
\end{document}
